% language=us

\environment context-2022-style

\startdocument
  [title={SQL in Metapost},
   banner={combining technologies},
   location={context\enspace {\bf 2022}\enspace meeting}]

\starttitle[title=The reason]

\startitemize

\startitem
    So far I had no real reason but there wat this question on the list about
    saving values.
\stopitem

\startitem
    A natural place to store values is in a \LUA\ table and one can easily save
    and load these with helpers.
\stopitem

\startitem
    Less obvious is to use a database, and a complication with for instance most
    \SQL\ databases, is that one has to have a service running.
\stopitem

\startitem
    A reasonable alternative is \SQLITE\ which uses a file and because it is part
    of the process it has some overhead.
\stopitem

\startitem
    Of course a database mostly makes sense with large quantities of data that
    needs to be accessed randomly using different lookup criteria.
\stopitem

\startitem
    The downside is that one needs to have the library available and but in this
    case it is a stable one and it does not influence rendering.
\stopitem

\stopitemize

\stoptitle

\starttitle[title=Libraries]

\startitemize

    \startitem
        When playing with this early in the \LUATEX\ project we started with
        native libraries that came with \LUA, and in particular \MYSQL. A simple
        command line client was used as a start.
    \stopitem

    \startitem
        The \LUA\ related libraries were sort of unreliable (and not really
        maintained) and not implemented that efficiently so I switched to just
        the main one that came with the client.
    \stopitem

    \startitem
        For a while we had \SWIGLIB\ which works okay but in order to use that an
        infrastructure is needed and we don't see that as our responsibility so
        that project sort of faded out.
    \stopitem

    \startitem
        In the end in \CONTEXT\ we used a minimal \FFI\ interface in \LUATEX.
    \stopitem

    \startitem
        In \LUAMETATEX\ we have a few so called optional libraries. These are
        in fact very lightweight interfaces that assume most work to be done in
        \LUA. We don't want dependencies and only these few are officially
        supported in this way. The \SQLITE\ interface is one of them.
    \stopitem

    \startitem
        With \CONTEXT\ we ship several \LUA\ support mechanisms for accessing
        databases. These perform well and we use it for web driven typesetting
        services (framework).
    \stopitem

\stopitemize

\stoptitle

\starttitle[title=Prerequisites]

We assume that the \type {sqlite3.dll} is present. Don't forget to remake the
file database after putting the file in place. On \MSWINDOWS\ the preferred
location is here:

\starttyping
E:/tex-context/tex/texmf-win64/bin/lib/luametatex/sqlite/sqlite3.dll
\stoptyping

If at some point loading fails one can turn on trackers:

\starttyping
\enabletrackers[sql*]
\enabletrackers[*lib*]
\stoptyping

As with any library, using one that comes from the system can be of any version
and who knows what happened to the \API.

\blank[3*big]

{\em Walk through the demo: context-2022-demo-sql-mps.tex}

\stoptitle

\stopdocument
